\scanpage{79}
La preuve de g) est \deleted{\unsure{évidemment|immédiatement}} facile (récurrence sur la \deleted{\unsure{nombre des}}\inserted{multiplicité} du pôle\deleted{\unsure{s}} à l'origine) \unsure{si on le sait|si on la sait|puisque la suite}, est au moins d'ordre 2).

On montre maintenant \inserted{aussi} facilement que \deleted{\unsure{toute}} fonction elliptique impaire est de la forme $\mathord{\unsure{\wp|p}}'R(\wp)$ -- $R$ fraction rationnelle -- (c'est évident d'après b) etc où $R$ est un polynôme si et seulement si l'ens des pôles est $\subset\unsure{\mathfrak P|\mathfrak B}$ (appliquer g))

\noindent \unsure{2)|II)|D)} \fbox{\uline{Compléments sur $\wp,\wp'$}}

\[
\left\{\begin{aligned}
\wp&=\frac1{z^2}+\sum_{m,n}'\left(\frac1{(z-2m\omega_1-2n\omega_2)^2}-\frac1{(2m\omega_1+2n\omega_2)^2}\right)\\
\wp'&=-2\sum_{m,n}\frac1{(z-2m\omega_1-2n\omega_2)^3}\deleted{\unsure{=-\frac2{z^3}+}}
\end{aligned}\right.\tag*{(1)}
\]
\[
\wp z=\frac1{z^2}\underset{\unsure{\times}}{+}c_1z^{\unsure{2|?}}+c_2z^4+\ldots\tag*{(2)}
\]
Les $c_i$ s'obtiennent \unsure{à|aisément à} partir de (1), en dérivant les deux expressions données de $\wp z$, et en comparant pour $z\to0$
\[
\left\{\begin{aligned}
c_1&=3\sum_{m,n}'\frac1{(2m\omega_1+2n\omega_2)^4}\\
&\quad -\quad -\quad -\\
c_k&=(2k+1)\sum_{m,n}'\frac1{(2m\omega_1+2n\omega_2)^{2\unsure{n|k}+2}}
\end{aligned}\right.\tag*{(3)}
\]
Les zéros de $\wp'(z)$ dans $P$ sont $\omega_1,\omega_2,-\omega_3=\omega_1+\omega_2$. On pose \fbox{$\wp\omega_i=e_i$} $\omega_i$ est racine double de l'équation $\wp z\deleted{\unsure{+}}-e_i=0$ ; il suit que \uline{les $e_i$ sont distincts}, \unsure{car|sinon} $\wp z-e_i$ n'a que $2$ racines (distinctes ou non à priori) dans $C/\mathfrak P$

% archive sheet number: 78
